% Options for packages loaded elsewhere
\PassOptionsToPackage{unicode}{hyperref}
\PassOptionsToPackage{hyphens}{url}
\documentclass[
]{article}
\usepackage{xcolor}
\usepackage[margin=1in]{geometry}
\usepackage{amsmath,amssymb}
\setcounter{secnumdepth}{5}
\usepackage{iftex}
\ifPDFTeX
  \usepackage[T1]{fontenc}
  \usepackage[utf8]{inputenc}
  \usepackage{textcomp} % provide euro and other symbols
\else % if luatex or xetex
  \usepackage{unicode-math} % this also loads fontspec
  \defaultfontfeatures{Scale=MatchLowercase}
  \defaultfontfeatures[\rmfamily]{Ligatures=TeX,Scale=1}
\fi
\usepackage{lmodern}
\ifPDFTeX\else
  % xetex/luatex font selection
\fi
% Use upquote if available, for straight quotes in verbatim environments
\IfFileExists{upquote.sty}{\usepackage{upquote}}{}
\IfFileExists{microtype.sty}{% use microtype if available
  \usepackage[]{microtype}
  \UseMicrotypeSet[protrusion]{basicmath} % disable protrusion for tt fonts
}{}
\makeatletter
\@ifundefined{KOMAClassName}{% if non-KOMA class
  \IfFileExists{parskip.sty}{%
    \usepackage{parskip}
  }{% else
    \setlength{\parindent}{0pt}
    \setlength{\parskip}{6pt plus 2pt minus 1pt}}
}{% if KOMA class
  \KOMAoptions{parskip=half}}
\makeatother
\usepackage{color}
\usepackage{fancyvrb}
\newcommand{\VerbBar}{|}
\newcommand{\VERB}{\Verb[commandchars=\\\{\}]}
\DefineVerbatimEnvironment{Highlighting}{Verbatim}{commandchars=\\\{\}}
% Add ',fontsize=\small' for more characters per line
\usepackage{framed}
\definecolor{shadecolor}{RGB}{248,248,248}
\newenvironment{Shaded}{\begin{snugshade}}{\end{snugshade}}
\newcommand{\AlertTok}[1]{\textcolor[rgb]{0.94,0.16,0.16}{#1}}
\newcommand{\AnnotationTok}[1]{\textcolor[rgb]{0.56,0.35,0.01}{\textbf{\textit{#1}}}}
\newcommand{\AttributeTok}[1]{\textcolor[rgb]{0.13,0.29,0.53}{#1}}
\newcommand{\BaseNTok}[1]{\textcolor[rgb]{0.00,0.00,0.81}{#1}}
\newcommand{\BuiltInTok}[1]{#1}
\newcommand{\CharTok}[1]{\textcolor[rgb]{0.31,0.60,0.02}{#1}}
\newcommand{\CommentTok}[1]{\textcolor[rgb]{0.56,0.35,0.01}{\textit{#1}}}
\newcommand{\CommentVarTok}[1]{\textcolor[rgb]{0.56,0.35,0.01}{\textbf{\textit{#1}}}}
\newcommand{\ConstantTok}[1]{\textcolor[rgb]{0.56,0.35,0.01}{#1}}
\newcommand{\ControlFlowTok}[1]{\textcolor[rgb]{0.13,0.29,0.53}{\textbf{#1}}}
\newcommand{\DataTypeTok}[1]{\textcolor[rgb]{0.13,0.29,0.53}{#1}}
\newcommand{\DecValTok}[1]{\textcolor[rgb]{0.00,0.00,0.81}{#1}}
\newcommand{\DocumentationTok}[1]{\textcolor[rgb]{0.56,0.35,0.01}{\textbf{\textit{#1}}}}
\newcommand{\ErrorTok}[1]{\textcolor[rgb]{0.64,0.00,0.00}{\textbf{#1}}}
\newcommand{\ExtensionTok}[1]{#1}
\newcommand{\FloatTok}[1]{\textcolor[rgb]{0.00,0.00,0.81}{#1}}
\newcommand{\FunctionTok}[1]{\textcolor[rgb]{0.13,0.29,0.53}{\textbf{#1}}}
\newcommand{\ImportTok}[1]{#1}
\newcommand{\InformationTok}[1]{\textcolor[rgb]{0.56,0.35,0.01}{\textbf{\textit{#1}}}}
\newcommand{\KeywordTok}[1]{\textcolor[rgb]{0.13,0.29,0.53}{\textbf{#1}}}
\newcommand{\NormalTok}[1]{#1}
\newcommand{\OperatorTok}[1]{\textcolor[rgb]{0.81,0.36,0.00}{\textbf{#1}}}
\newcommand{\OtherTok}[1]{\textcolor[rgb]{0.56,0.35,0.01}{#1}}
\newcommand{\PreprocessorTok}[1]{\textcolor[rgb]{0.56,0.35,0.01}{\textit{#1}}}
\newcommand{\RegionMarkerTok}[1]{#1}
\newcommand{\SpecialCharTok}[1]{\textcolor[rgb]{0.81,0.36,0.00}{\textbf{#1}}}
\newcommand{\SpecialStringTok}[1]{\textcolor[rgb]{0.31,0.60,0.02}{#1}}
\newcommand{\StringTok}[1]{\textcolor[rgb]{0.31,0.60,0.02}{#1}}
\newcommand{\VariableTok}[1]{\textcolor[rgb]{0.00,0.00,0.00}{#1}}
\newcommand{\VerbatimStringTok}[1]{\textcolor[rgb]{0.31,0.60,0.02}{#1}}
\newcommand{\WarningTok}[1]{\textcolor[rgb]{0.56,0.35,0.01}{\textbf{\textit{#1}}}}
\usepackage{graphicx}
\makeatletter
\newsavebox\pandoc@box
\newcommand*\pandocbounded[1]{% scales image to fit in text height/width
  \sbox\pandoc@box{#1}%
  \Gscale@div\@tempa{\textheight}{\dimexpr\ht\pandoc@box+\dp\pandoc@box\relax}%
  \Gscale@div\@tempb{\linewidth}{\wd\pandoc@box}%
  \ifdim\@tempb\p@<\@tempa\p@\let\@tempa\@tempb\fi% select the smaller of both
  \ifdim\@tempa\p@<\p@\scalebox{\@tempa}{\usebox\pandoc@box}%
  \else\usebox{\pandoc@box}%
  \fi%
}
% Set default figure placement to htbp
\def\fps@figure{htbp}
\makeatother
\setlength{\emergencystretch}{3em} % prevent overfull lines
\providecommand{\tightlist}{%
  \setlength{\itemsep}{0pt}\setlength{\parskip}{0pt}}
\usepackage{bookmark}
\IfFileExists{xurl.sty}{\usepackage{xurl}}{} % add URL line breaks if available
\urlstyle{same}
\hypersetup{
  pdftitle={Using Keeper with Large Language Models},
  pdfauthor={Martijn Schuemie},
  hidelinks,
  pdfcreator={LaTeX via pandoc}}

\title{Using Keeper with Large Language Models}
\author{Martijn Schuemie}
\date{2026-03-26}

\begin{document}
\maketitle

{
\setcounter{tocdepth}{2}
\tableofcontents
}
\section{Introduction}\label{introduction}

This vignette demonstrates how to use \textbf{Keeper} to generate
patient summaries and systematically review them using a Large Language
Model (LLM).

\textbf{Prerequisites:} Before proceeding, ensure you have read the
\emph{`Generating Keeper profiles for review'} vignette. We will build
directly on the Type 1 Diabetes Mellitus (T1DM) example established
there, assuming you have already created the necessary objects and input
concept sets.

\begin{center}\rule{0.5\linewidth}{0.5pt}\end{center}

\section{Workflow 1: Direct Cohort
Review}\label{workflow-1-direct-cohort-review}

In this first workflow, we will take an existing cohort, generate Keeper
profiles, and use an LLM to evaluate those specific cases. We assume
that Keeper has already been executed to produce the following output:

\begin{Shaded}
\begin{Highlighting}[]
\NormalTok{keeper}
\end{Highlighting}
\end{Shaded}

\begin{verbatim}
## # A tibble: 2,234 x 8
##    generatedId startDay endDay conceptId conceptName category target              extraData
##          <dbl>    <dbl>  <dbl>     <dbl> <chr>       <chr>    <chr>               <chr>    
##  1          16        0      0        NA 67          age      Disease of interest <NA>     
##  2           8        0      0        NA 70          age      Disease of interest <NA>     
##  3           6        0      0        NA 72          age      Disease of interest <NA>     
##  4          14        0      0        NA 61          age      Disease of interest <NA>     
##  5          10        0      0        NA 68          age      Disease of interest <NA>     
##  6           5        0      0        NA 78          age      Disease of interest <NA>     
##  7           7        0      0        NA 70          age      Disease of interest <NA>     
##  8          15        0      0        NA 81          age      Disease of interest <NA>     
##  9           9        0      0        NA 79          age      Disease of interest <NA>     
## 10          13        0      0        NA 79          age      Disease of interest <NA>     
## # i 2,224 more rows
\end{verbatim}

\subsection{1. Choosing an LLM Provider}\label{choosing-an-llm-provider}

Keeper integrates with the \texttt{ellmer} package, allowing you to
connect to the LLM provider of your choice (e.g., Anthropic, Google,
OpenAI, or a local model).

For instance, to connect to OpenAI's ChatGPT, you can use:

\begin{Shaded}
\begin{Highlighting}[]
\FunctionTok{library}\NormalTok{(ellmer)}
\NormalTok{client }\OtherTok{\textless{}{-}} \FunctionTok{chat\_openai}\NormalTok{()}
\end{Highlighting}
\end{Shaded}

\emph{Note:} This requires the \texttt{OPENAI\_API\_KEY} environmental
variable to be set. For detailed connection instructions across various
providers, refer to the \href{https://ellmer.tidyverse.org}{ellmer
package documentation}.

\textbf{Model Requirements:} The selected LLM must be capable of
advanced reasoning and producing structured outputs. In our testing,
OpenAI's GPT-o3 yielded the best results.

\subsection{2. Executing the Review}\label{executing-the-review}

Once connected, we can command the LLM to review the patient cases:

\begin{Shaded}
\begin{Highlighting}[]
\NormalTok{promptSettings }\OtherTok{\textless{}{-}} \FunctionTok{createPromptSettings}\NormalTok{()}
\NormalTok{llmReviews }\OtherTok{\textless{}{-}} \FunctionTok{reviewCases}\NormalTok{(}
  \AttributeTok{keeper =}\NormalTok{ keeper,}
  \AttributeTok{settings =}\NormalTok{ promptSettings,}
  \AttributeTok{phenotypeName =} \StringTok{"Type I Diabetes Mellitus (T1DM)"}\NormalTok{,}
  \AttributeTok{client =}\NormalTok{ client,}
  \AttributeTok{cacheFolder =} \StringTok{"cache"}
\NormalTok{)}
\end{Highlighting}
\end{Shaded}

\begin{verbatim}
## Reviewing person 1 of 20
## ...
## Reviewing person 20 of 20
## Reviewing cases took 3.9 mins and cost $0.50
\end{verbatim}

\textbf{Key Parameters:} * \texttt{phenotypeName}: This string is
inserted directly into the LLM prompt. It must exactly match the disease
you are adjudicating. * \texttt{cacheFolder}: A local directory where
prompts and responses are safely cached. If the process is interrupted,
re-running the function will skip previously generated responses.

We can inspect the resulting adjudications:

\begin{Shaded}
\begin{Highlighting}[]
\NormalTok{llmReviews}
\end{Highlighting}
\end{Shaded}

\begin{verbatim}
## # A tibble: 20 x 9
##    generatedId phenotype                       isCase certainty indexDay justification                                                                                                           cohortPrevalence model keeperVersion
##          <dbl> <chr>                           <chr>  <chr>        <dbl> <chr>                                                                                                                              <dbl> <chr> <chr>        
##  1          16 Type I Diabetes Mellitus (T1DM) yes    high             0 "The patient is a 67-year-old male.  Evidence FOR true diabetes:  1. Multiple, longitudinal diagnostic entries explici~           0.0488 o3    2.0.0        
##  2           8 Type I Diabetes Mellitus (T1DM) no     high            NA "Evidence suggesting T1DM:\n• ICD codes for “Type 1 diabetes mellitus” appear on day 0 (hospital stay) and day 93.\n• ~           0.0488 o3    2.0.0        
##  3           6 Type I Diabetes Mellitus (T1DM) no     high            NA "Only a single code for “Type 1 diabetes mellitus” appears (secondary diagnosis on one outpatient claim). The patient ~           0.0488 o3    2.0.0        
##  4          14 Type I Diabetes Mellitus (T1DM) no     high            NA "• Only one code for Type 1 DM (day 0) and it is listed as a secondary diagnosis. A single occurrence—especially as a ~           0.0488 o3    2.0.0        
##  5          10 Type I Diabetes Mellitus (T1DM) no     high            NA "The patient is a 68-year-old man whose diabetes was first treated (≥294 days before index) with classic type 2 agents~           0.0488 o3    2.0.0        
##  6           5 Type I Diabetes Mellitus (T1DM) yes    low              0 "• The patient carries the diagnosis code “Type 1 diabetes mellitus” on six separate service dates spread over >4 year~           0.0488 o3    2.0.0        
##  7           7 Type I Diabetes Mellitus (T1DM) no     high            NA "Only a single mention of “Type 1 diabetes mellitus” appears on day 0 and it is listed as a secondary (not primary) di~           0.0488 o3    2.0.0        
##  8          15 Type I Diabetes Mellitus (T1DM) yes    high             0 "The patient clearly has diabetes (multiple HbA1c tests, diabetic kidney disease later on).  The key question is the T~           0.0488 o3    2.0.0        
##  9           9 Type I Diabetes Mellitus (T1DM) no     high            NA "• Only a single encounter (day 0 ophthalmology outpatient visit) contains a Type 1 diabetes mellitus (T1DM) code, and~           0.0488 o3    2.0.0        
## 10          13 Type I Diabetes Mellitus (T1DM) no     high            NA "Evidence against true Type 1 DM: 1) Patient is 79 y male—de novo T1DM onset at this age is exceedingly rare; most old~           0.0488 o3    2.0.0        
## 11          17 Type I Diabetes Mellitus (T1DM) no     high            NA "Evidence in favor of T1DM: (1) Single hospitalization (days 0-11) lists secondary codes for “Type 1 diabetes mellitus~           0.0488 o3    2.0.0        
## 12          18 Type I Diabetes Mellitus (T1DM) no     high            NA "Only a single outpatient secondary diagnosis code for Type 1 diabetes appears (day 0).  • Age 81 makes new-onset T1DM~           0.0488 o3    2.0.0        
## 13          19 Type I Diabetes Mellitus (T1DM) no     low             NA "There are only two encounters over a 10-month span that carry a Type 1 diabetes (T1DM) diagnosis code, one as a secon~           0.0488 o3    2.0.0        
## 14           3 Type I Diabetes Mellitus (T1DM) no     high            NA "Multiple Type-1 DM diagnosis codes are present over several years and the patient does receive prolonged insulin ther~           0.0488 o3    2.0.0        
## 15           1 Type I Diabetes Mellitus (T1DM) yes    high             0 "Multiple independent pieces of evidence support true, ongoing Type 1 diabetes rather than mere rule-out testing. 1) R~           0.0488 o3    2.0.0        
## 16           2 Type I Diabetes Mellitus (T1DM) no     high            NA "Multiple T1DM diagnosis codes appear, but the therapeutic and coding pattern fits long-standing Type 2 diabetes, not ~           0.0488 o3    2.0.0        
## 17          12 Type I Diabetes Mellitus (T1DM) no     high            NA "The patient is a 70-year-old woman with a long history of oral antihyperglycaemic drugs that are standard for Type 2 ~           0.0488 o3    2.0.0        
## 18          11 Type I Diabetes Mellitus (T1DM) no     high            NA "Evidence FOR true T1DM: (1) Diagnosis code for “Type 1 diabetes mellitus” appears on index visit (day 0) and twice la~           0.0488 o3    2.0.0        
## 19          20 Type I Diabetes Mellitus (T1DM) yes    high             0 "The patient is a 72-year-old woman who is treated with insulin continuously throughout the entire observation window.~           0.0488 o3    2.0.0        
## 20           4 Type I Diabetes Mellitus (T1DM) no     high            NA "Only a single mention of “Type 1 diabetes mellitus” appears on day 0, and it is listed merely as a secondary outpatie~           0.0488 o3    2.0.0
\end{verbatim}

With this data, computing the positive predictive value (PPV) of the
reviewed cohort is straightforward:

\begin{Shaded}
\begin{Highlighting}[]
\FunctionTok{mean}\NormalTok{(llmReviews}\SpecialCharTok{$}\NormalTok{isCase }\SpecialCharTok{==} \StringTok{"yes"}\NormalTok{)}
\end{Highlighting}
\end{Shaded}

\begin{verbatim}
## [1] 0.25
\end{verbatim}

\begin{center}\rule{0.5\linewidth}{0.5pt}\end{center}

\section{Workflow 2: Validating with a Reference
Cohort}\label{workflow-2-validating-with-a-reference-cohort}

Instead of analyzing an existing cohort just to find its PPV, we can
build a \textbf{Highly-Sensitive Cohort (HSC)}. An HSC is designed to be
overly broad, ensuring no true cases are missed.

By taking a large sample of this HSC and having an LLM review it, we
create a gold-standard \textbf{Reference Cohort}. This reference cohort
can then be used to rapidly evaluate the PPV \emph{and} sensitivity of
any other cohort definitions for the same phenotype.

\begin{figure}
\centering
\includegraphics[width=3.64583in,height=\textheight,keepaspectratio]{Diagram.png}
\caption{Using a reference cohort (a sample of a highly-sensitive
cohort) to evaluate cohorts for a phenotype.}
\end{figure}

\subsection{1. Creating the Highly-Sensitive
Cohort}\label{creating-the-highly-sensitive-cohort}

The \texttt{createSensitiveCohort()} function automatically constructs
this broad cohort using the same input concept sets utilized by
\texttt{generateKeeper()}.

The cohort is constructed by combining two sets:

\begin{enumerate}
\def\labelenumi{\arabic{enumi}.}
\item
  Anybody who has one of the condition concepts of the disease of
  interest, indexed on the first occurrence.
\item
  Anybody who has `highly specific concepts' from at least two Keeper
  categories (symptoms, drugs, etc.), indexed on the first occurrence.
\end{enumerate}

Highly specific concepts are concepts where at least 10\% of those who
have that concept also have a condition concept for the disease of
interest.

\begin{Shaded}
\begin{Highlighting}[]
\FunctionTok{createSensitiveCohort}\NormalTok{(}
  \AttributeTok{connectionDetails =}\NormalTok{ connectionDetails,}
  \AttributeTok{cdmDatabaseSchema =}\NormalTok{ cdmDatabaseSchema,}
  \AttributeTok{cohortDatabaseSchema =}\NormalTok{ cohortDatabaseSchema,}
  \AttributeTok{cohortTable =}\NormalTok{ cohortTable,}
  \AttributeTok{cohortDefinitionId =} \DecValTok{2}\NormalTok{,}
  \AttributeTok{createCohortTable =} \ConstantTok{FALSE}\NormalTok{,}
  \AttributeTok{keeperConceptSets =}\NormalTok{ conceptSets}
\NormalTok{)}
\end{Highlighting}
\end{Shaded}

Here, we assign \texttt{cohortDefinitionId\ =\ 2} to prevent overwriting
the simpler cohort (ID 1) we created in the previous vignette.

\subsection{2. Running Keeper on the
HSC}\label{running-keeper-on-the-hsc}

Next, we extract the patient profiles for our new sensitive cohort:

\begin{Shaded}
\begin{Highlighting}[]
\NormalTok{keeperHsc }\OtherTok{\textless{}{-}} \FunctionTok{generateKeeper}\NormalTok{(}
  \AttributeTok{connectionDetails =}\NormalTok{ connectionDetails,}
  \AttributeTok{cohortDatabaseSchema =}\NormalTok{ cohortDatabaseSchema,}
  \AttributeTok{cdmDatabaseSchema =}\NormalTok{ cdmDatabaseSchema,}
  \AttributeTok{cohortTable =}\NormalTok{ cohortTable,}
  \AttributeTok{cohortDefinitionId =} \DecValTok{2}\NormalTok{,}
  \AttributeTok{sampleSize =} \DecValTok{10000}\NormalTok{,}
  \AttributeTok{phenotypeName =} \StringTok{"Type I Diabetes Mellitus (T1DM)"}\NormalTok{,}
  \AttributeTok{keeperConceptSets =}\NormalTok{ conceptSets,}
  \AttributeTok{removePii =} \ConstantTok{FALSE}
\NormalTok{)}
\end{Highlighting}
\end{Shaded}

\textbf{Important considerations:} * \textbf{Sample Size:} We sample
10,000 persons to guarantee enough true positive cases for robust PPV
and sensitivity calculations. * \textbf{PII Handling:} We set
\texttt{removePii\ =\ FALSE} so that the resulting LLM adjudications can
be joined accurately with specific target cohorts later during
evaluation.

\subsection{3. Adjudicating the HSC via
LLM}\label{adjudicating-the-hsc-via-llm}

We now send all 10,000 sampled profiles to the LLM for review.

\begin{quote}
\textbf{Cost Warning:} Adjudicating a dataset of this size can incur
significant API costs. For example, using our preferred model (GPT-o3)
to review 10,000 cases typically costs between \$60 and \$80.
\end{quote}

\begin{Shaded}
\begin{Highlighting}[]
\NormalTok{promptSettings }\OtherTok{\textless{}{-}} \FunctionTok{createPromptSettings}\NormalTok{()}
\NormalTok{llmReviewsHsc }\OtherTok{\textless{}{-}} \FunctionTok{reviewCases}\NormalTok{(}
  \AttributeTok{keeper =}\NormalTok{ keeperHsc,}
  \AttributeTok{settings =}\NormalTok{ promptSettings,}
  \AttributeTok{client =}\NormalTok{ client,}
  \AttributeTok{cacheFolder =} \StringTok{"cacheVignetteHsc"}
\NormalTok{)}
\end{Highlighting}
\end{Shaded}

Once complete, this dataset becomes our \emph{Reference Cohort}.

\subsection{4. Storing the Reference
Cohort}\label{storing-the-reference-cohort}

To use the reference cohort across different analyses, we upload it back
to our database server:

\begin{Shaded}
\begin{Highlighting}[]
\NormalTok{referenceCohortDatabaseSchema }\OtherTok{\textless{}{-}} \StringTok{"cdm"}
\NormalTok{referenceCohortTable }\OtherTok{\textless{}{-}} \StringTok{"reference\_cohort"}

\FunctionTok{uploadReferenceCohort}\NormalTok{(}
  \AttributeTok{connectionDetails =}\NormalTok{ connectionDetails,}
  \AttributeTok{referenceCohortDatabaseSchema =}\NormalTok{ referenceCohortDatabaseSchema,}
  \AttributeTok{referenceCohortTableNames =} \FunctionTok{createReferenceCohortTableNames}\NormalTok{(referenceCohortTable),}
  \AttributeTok{referenceCohortDefinitionId =} \DecValTok{1}\NormalTok{,}
  \AttributeTok{createReferenceCohortTables =} \ConstantTok{TRUE}\NormalTok{,}
  \AttributeTok{reviews =}\NormalTok{ llmReviewsHsc}
\NormalTok{)}
\end{Highlighting}
\end{Shaded}

Setting \texttt{createReferenceCohortTables\ =\ TRUE} generates two
tables (\texttt{reference\_cohort} and
\texttt{reference\_cohort\_metadata}). We also assign a
\texttt{referenceCohortDefinitionId} of \texttt{1} to uniquely identify
it for future comparisons.

\subsection{5. Evaluating Target
Cohorts}\label{evaluating-target-cohorts}

With our reference cohort safely stored in the database, evaluating any
standard T1DM cohort is highly efficient. Let's evaluate the simple
target cohort (\texttt{cohortDefinitionId\ =\ 1}) from the previous
vignette:

\begin{Shaded}
\begin{Highlighting}[]
\NormalTok{metrics }\OtherTok{\textless{}{-}} \FunctionTok{computeCohortOperatingCharacteristics}\NormalTok{(}
  \AttributeTok{connectionDetails =}\NormalTok{ connectionDetails,}
  \AttributeTok{cohortDatabaseSchema =}\NormalTok{ cohortDatabaseSchema,}
  \AttributeTok{cohortTable =}\NormalTok{ cohortTable,}
  \AttributeTok{cohortDefinitionId =} \DecValTok{1}\NormalTok{,}
  \AttributeTok{referenceCohortDatabaseSchema =}\NormalTok{ referenceCohortDatabaseSchema,}
  \AttributeTok{referenceCohortTableNames =} \FunctionTok{createReferenceCohortTableNames}\NormalTok{(referenceCohortTable),}
  \AttributeTok{referenceCohortDefinitionId =} \DecValTok{1}
\NormalTok{)}
\end{Highlighting}
\end{Shaded}

Finally, we can view the computed performance metrics (such as
sensitivity, specificity, and PPV) for our target cohort:

\begin{Shaded}
\begin{Highlighting}[]
\NormalTok{metrics}
\end{Highlighting}
\end{Shaded}

\begin{verbatim}
## # A tibble: 3 x 21
##   certainty    tp    fp    tn    fn sensitivity sensitivityLb sensitivityUb specificity specificityLb specificityUb specificityOverall specificityOverallLb specificityOverallUb   ppv ppvLb ppvUb   auc kappa prevalence
##   <chr>     <dbl> <dbl> <dbl> <dbl>       <dbl>         <dbl>         <dbl>       <dbl>         <dbl>         <dbl>              <dbl>                <dbl>                <dbl> <dbl> <dbl> <dbl> <dbl> <dbl>      <dbl>
## 1 low         100    20    57    55       0.645         0.564         0.720       0.740         0.628         0.834              0.910                0.865                0.944 0.833 0.754 0.895 0.693 0.346      0.672
## 2 high        941   215  8358     7       0.993         0.985         0.997       0.975         0.971         0.978              0.993                0.992                0.994 0.814 0.790 0.836 0.984 0.882      0.107
## 3 all        1041   235  8415    62       0.944         0.929         0.957       0.973         0.969         0.976              0.992                0.991                0.993 0.816 0.793 0.837 0.958 0.858      0.121
## # i 1 more variable: prevalenceOverall <dbl>
\end{verbatim}

\end{document}
